% Options for packages loaded elsewhere
% Options for packages loaded elsewhere
\PassOptionsToPackage{unicode}{hyperref}
\PassOptionsToPackage{hyphens}{url}
\PassOptionsToPackage{dvipsnames,svgnames,x11names}{xcolor}
%
\documentclass[
]{article}
\usepackage{xcolor}
\usepackage{amsmath,amssymb}
\setcounter{secnumdepth}{-\maxdimen} % remove section numbering
\usepackage{iftex}
\ifPDFTeX
  \usepackage[T1]{fontenc}
  \usepackage[utf8]{inputenc}
  \usepackage{textcomp} % provide euro and other symbols
\else % if luatex or xetex
  \usepackage{unicode-math} % this also loads fontspec
  \defaultfontfeatures{Scale=MatchLowercase}
  \defaultfontfeatures[\rmfamily]{Ligatures=TeX,Scale=1}
\fi
\usepackage{lmodern}
\ifPDFTeX\else
  % xetex/luatex font selection
\fi
% Use upquote if available, for straight quotes in verbatim environments
\IfFileExists{upquote.sty}{\usepackage{upquote}}{}
\IfFileExists{microtype.sty}{% use microtype if available
  \usepackage[]{microtype}
  \UseMicrotypeSet[protrusion]{basicmath} % disable protrusion for tt fonts
}{}
\makeatletter
\@ifundefined{KOMAClassName}{% if non-KOMA class
  \IfFileExists{parskip.sty}{%
    \usepackage{parskip}
  }{% else
    \setlength{\parindent}{0pt}
    \setlength{\parskip}{6pt plus 2pt minus 1pt}}
}{% if KOMA class
  \KOMAoptions{parskip=half}}
\makeatother
% Make \paragraph and \subparagraph free-standing
\makeatletter
\ifx\paragraph\undefined\else
  \let\oldparagraph\paragraph
  \renewcommand{\paragraph}{
    \@ifstar
      \xxxParagraphStar
      \xxxParagraphNoStar
  }
  \newcommand{\xxxParagraphStar}[1]{\oldparagraph*{#1}\mbox{}}
  \newcommand{\xxxParagraphNoStar}[1]{\oldparagraph{#1}\mbox{}}
\fi
\ifx\subparagraph\undefined\else
  \let\oldsubparagraph\subparagraph
  \renewcommand{\subparagraph}{
    \@ifstar
      \xxxSubParagraphStar
      \xxxSubParagraphNoStar
  }
  \newcommand{\xxxSubParagraphStar}[1]{\oldsubparagraph*{#1}\mbox{}}
  \newcommand{\xxxSubParagraphNoStar}[1]{\oldsubparagraph{#1}\mbox{}}
\fi
\makeatother


\usepackage{longtable,booktabs,array}
\usepackage{calc} % for calculating minipage widths
% Correct order of tables after \paragraph or \subparagraph
\usepackage{etoolbox}
\makeatletter
\patchcmd\longtable{\par}{\if@noskipsec\mbox{}\fi\par}{}{}
\makeatother
% Allow footnotes in longtable head/foot
\IfFileExists{footnotehyper.sty}{\usepackage{footnotehyper}}{\usepackage{footnote}}
\makesavenoteenv{longtable}
\usepackage{graphicx}
\makeatletter
\newsavebox\pandoc@box
\newcommand*\pandocbounded[1]{% scales image to fit in text height/width
  \sbox\pandoc@box{#1}%
  \Gscale@div\@tempa{\textheight}{\dimexpr\ht\pandoc@box+\dp\pandoc@box\relax}%
  \Gscale@div\@tempb{\linewidth}{\wd\pandoc@box}%
  \ifdim\@tempb\p@<\@tempa\p@\let\@tempa\@tempb\fi% select the smaller of both
  \ifdim\@tempa\p@<\p@\scalebox{\@tempa}{\usebox\pandoc@box}%
  \else\usebox{\pandoc@box}%
  \fi%
}
% Set default figure placement to htbp
\def\fps@figure{htbp}
\makeatother





\setlength{\emergencystretch}{3em} % prevent overfull lines

\providecommand{\tightlist}{%
  \setlength{\itemsep}{0pt}\setlength{\parskip}{0pt}}



 


\makeatletter
\@ifpackageloaded{caption}{}{\usepackage{caption}}
\AtBeginDocument{%
\ifdefined\contentsname
  \renewcommand*\contentsname{Table of contents}
\else
  \newcommand\contentsname{Table of contents}
\fi
\ifdefined\listfigurename
  \renewcommand*\listfigurename{List of Figures}
\else
  \newcommand\listfigurename{List of Figures}
\fi
\ifdefined\listtablename
  \renewcommand*\listtablename{List of Tables}
\else
  \newcommand\listtablename{List of Tables}
\fi
\ifdefined\figurename
  \renewcommand*\figurename{Figure}
\else
  \newcommand\figurename{Figure}
\fi
\ifdefined\tablename
  \renewcommand*\tablename{Table}
\else
  \newcommand\tablename{Table}
\fi
}
\@ifpackageloaded{float}{}{\usepackage{float}}
\floatstyle{ruled}
\@ifundefined{c@chapter}{\newfloat{codelisting}{h}{lop}}{\newfloat{codelisting}{h}{lop}[chapter]}
\floatname{codelisting}{Listing}
\newcommand*\listoflistings{\listof{codelisting}{List of Listings}}
\makeatother
\makeatletter
\makeatother
\makeatletter
\@ifpackageloaded{caption}{}{\usepackage{caption}}
\@ifpackageloaded{subcaption}{}{\usepackage{subcaption}}
\makeatother
\usepackage{bookmark}
\IfFileExists{xurl.sty}{\usepackage{xurl}}{} % add URL line breaks if available
\urlstyle{same}
\hypersetup{
  pdftitle={Analisis Rasch Menggunakna eRm},
  pdfauthor={The Curious Mind},
  colorlinks=true,
  linkcolor={blue},
  filecolor={Maroon},
  citecolor={Blue},
  urlcolor={Blue},
  pdfcreator={LaTeX via pandoc}}


\title{Analisis Rasch Menggunakna eRm}
\author{The Curious Mind}
\date{2026-06-10}
\begin{document}
\maketitle


\subsubsection{1. Andersen's Likelihood Ratio
Test}\label{andersens-likelihood-ratio-test}

Uji ini dilakukan untuk memvalidasi pemenuhan asumsi invariansi secara
empiris dalam kerangka \emph{Conditional Maximum Likelihood} (CML),
\textbf{Andersen's Likelihood Ratio (LR) Test} digunakan sebagai
instrumen diagnostik global. Uji ini mengevaluasi stabilitas parameter
soal dengan mempartisi sampel total ke dalam dua atau lebih sub-kelompok
berdasarkan kriteria tertentu (seperti skor median, mean atau variabel
demografis).

Secara matematis, prosedur ini membandingkan fungsi \emph{likelihood}
dari model yang diestimasi secara terpisah pada masing-masing
sub-kelompok dengan model gabungan. Apabila instrumen pengukuran valid
dan tidak bias, parameter kesulitan butir soal tidak akan menunjukkan
perbedaan yang signifikan secara statistik di antara sub-kelompok
tersebut, yang dibuktikan melalui nilai uji Chi-Square (\(\chi^2\))
dengan non-signifikansi (\(p > 0.05\)).

\begin{verbatim}

Andersen LR-test: 
LR-value: 10.051 
Chi-square df: 11 
p-value:  0.526 
\end{verbatim}

Meskipun nilai-p sebesar 0,526 secara konvensional mengindikasikan
kesesuaian model yang sempurna, algoritma CML menghasilkan peringatan
struktural yang kritis: sembilan item (S1, S2, S3, S4, S5, S7, S8, S9,
dan S12) secara mekanis dikeluarkan dari perhitungan karena pola respons
yang tidak sesuai dalam subkelompok yang terpartisi.

Ini terjadi karena paket \texttt{eRm} menggunakan metode estimasi
\textbf{Conditional Maximum Likelihood (CML)}. CML adalah mesin
matematika yang sangat kaku dan menuntut kesempurnaan distribusi data
pada setiap jenjang jawaban (karena instrumen tes menggunakan
\emph{Partial Credit Model} yang bertingkat).

Ketika model membagi siswa menjadi dua kelompok berdasarkan \textbf{skor
rata-rata (mean)}, yaitu kelompok siswa berkemampuan di bawah rata-rata
dan kelompok siswa di atas rata-rata, tercipta polarisasi yang tajam.
Pada 9 soal tersebut, terjadi \textbf{category sparsity (kekosongan sel
respons)} pada sub-kelompok tertentu:

\begin{itemize}
\item
  Pada soal yang sangat menantang, \emph{sama sekali tidak ada} siswa di
  kelompok bawah yang berhasil mencapai jenjang skor tertinggi (misal:
  skor 3 atau 4).
\item
  Sebaliknya, pada soal yang sangat mudah, \emph{sama sekali tidak ada}
  siswa di kelompok atas yang mendapatkan skor terendah (skor 0).
\end{itemize}

Dalam matematika CML, jika sebuah kategori jawaban memiliki frekuensi
nol di salah satu sub-kelompok, komputer tidak dapat menghitung titik
potong (threshold) karena akan membagi dengan angka nol. Agar sistem
tidak \emph{crash}, \texttt{eRm} secara otomatis menyingkirkan ke-9 soal
tersebut dan hanya menyisakan soal S6, S10, dan S11 yang datanya
``penuh'' di kedua kelompok.

Karena tujuan riset ini bukan untuk menstandarisasi tes (di mana kita
memaksakan semua soal harus seragam dan membuang yang meleset),
melainkan untuk \textbf{memetakan realitas interaksi antara siswa dan
kualitas soal}, anomali ini justru menjadi temuan utama.

Kita dapat mengeksplorasi jurang pemisah antara kualitas soal yang
dirancang oleh pembuat soal ahli dan pola pikir atau kemampuan faktual
dari sampel siswa.

\subsubsection{2. Uji Unidimensionality}\label{uji-unidimensionality}

Salah satu asumsi paling fundamental yang harus dipenuhi adalah
\textbf{unidimensionalitas}. Asumsi ini menyatakan bahwa seperangkat
butir soal (item) yang dianalisis harus mengukur \textbf{satu konstruk
atau sifat laten tunggal yang sama} (misalnya, hanya mengukur kemampuan
penalaran logis, bukan bercampur dengan kemampuan membaca atau kecemasan
numerik). Jika instrumen mengukur lebih dari satu dimensi secara
bersamaan, model matematika penskalaan tunggal akan kehilangan validitas
strukturalnya.

Untuk menguji asumsi ini secara empiris, \textbf{Martin-Löf Test}
digunakan sebagai instrumen diagnostik yang sangat andal di dalam
ekosistem \emph{Rasch Measurement}. Uji ini bekerja dengan membagi
perangkat soal (item) menjadi dua sub-kelompok (misalnya, membagi item
berdasarkan tingkat kesulitan menjadi kelompok soal mudah dan soal
sulit, atau membagi berdasarkan karakteristik konten secara logis).

Komputer kemudian akan mengestimasi parameter kemampuan responden
(\(\theta\)) secara terpisah menggunakan kedua sub-kelompok soal
tersebut. Logika dasarnya adalah: jika instrumen benar-benar
unidimensional, kemampuan seorang siswa yang diukur menggunakan separuh
soal pertama harus memprediksi secara presisi kemampuannya saat diukur
menggunakan separuh soal kedua.

Secara statistik, pengujian ini membandingkan \emph{likelihood} model
multidimensional dengan model unidimensional:

\begin{itemize}
\item
  \textbf{Hipotesis Nol (}\(H_0\)): Instrumen bersifat unidimensional
  (homogen).
\item
  \textbf{Hipotesis Alternatif (}\(H_1\)): Instrumen bersifat
  multidimensional (terdapat dimensi lain yang tidak terdeteksi).
\end{itemize}

Berbeda dari uji signifikansi pada umumnya di mana kita mencari nilai
\(p < 0.05\), \textbf{pada uji Martin-Löf kita menginginkan hasil yang
tidak signifikan secara statistik (}\(p > 0.05\)). Nilai \(p\) yang
besar menunjukkan bahwa kita gagal menolak \(H_0\), yang berarti
instrumen telah memenuhi asumsi kemurnian pengukuran
(unidimensionalitas).

\begin{verbatim}

Martin-Loef-Test (split criterion: median)
LR-value: 166.622 
Chi-square df: 527 
p-value: 1 
\end{verbatim}

\begin{itemize}
\item
  Dalam konteks uji kecocokan \emph{Martin-Löf}, nilai \texttt{p:\ 1}
  memberikan konfirmasi absolut bahwa \textbf{tidak ada sama sekali
  pelanggaran terhadap asumsi unidimensionalitas}. Data respons siswa
  tidak menunjukkan deviasi atau jurang pemisah dimensi yang menyimpang
  dari model teoretis Rasch.
\item
  Nilai \(LR\) jauh lebih kecil daripada nilai \(df\)
  (\(166.622 \ll 527\)) mengindikasikan bahwa varians sisa
  (\emph{residual variance}) di antara kedua sub-kelompok soal jauh
  lebih sempit atau lebih terkontrol daripada yang diprediksi oleh batas
  fluktuasi sampel acak. Hal ini memperkuat bukti bahwa kedua bagian
  instrumen mengukur konstruk yang identik tanpa ada gangguan dimensi
  sekunder.
\item
  Terlepas dari adanya variasi respons dan \emph{category sparsity} yang
  sempat terdeteksi pada kalkulasi parameter CML sub-kelompok, ke-12
  butir soal yang dirancang terbukti secara empiris mengukur satu
  dimensi laten yang sama persis.
\item
  Instrumen ini tidak mengalami \emph{multidimensionality leakage}.
  Hasil penskalaan tingkat kesulitan butir soal (\emph{item difficulty})
  dan estimasi kemampuan siswa (\emph{person ability}) berada pada satu
  sumbu kontinum logit yang murni, sehingga analisis lanjutan seperti
  \emph{Person-Item Map} (Wright Map) dapat diinterpretasikan.
\end{itemize}

\subsubsection{3. Item fit}\label{item-fit}

Setelah memastikan instrumen memenuhi asumsi kemurnian pengukuran
(unidimensionalitas) melalui uji Martin-Löf, langkah berikutnya dalam
pemodelan Rasch adalah mengevaluasi \textbf{kecocokan butir soal
(\emph{item fit})}.

\paragraph{3.1 Uji Validitas}\label{uji-validitas}

Validitas struktural dan keandalan instrumen dievaluasi lebih lanjut
melalui analisis kecocokan model (\emph{item-person fit}) menggunakan
pendekatan \emph{Partial Credit Model} (PCM). Evaluasi ini didasarkan
pada statistik \emph{Infit} dan \emph{Outfit Mean Square} (MSQ) serta
indeks standardisasi (\emph{t-fit} atau \emph{zstd}). Dalam kerangka
kerja eksploratif, analisis ini tidak ditujukan semata-mata untuk
reduksi instrumen secara kaku, melainkan untuk memetakan dinamika
interaksi antara kualitas soal rancangan dan pola pikir faktual
responden (\(N = 102\)).

\begin{verbatim}

Itemfit Statistics: 
      Chisq  df p-value Outfit MSQ Infit MSQ Outfit t Infit t Discrim
S1   76.388 101   0.968      0.749     0.785   -1.909  -2.282   0.606
S2   93.314 101   0.694      0.915     0.962   -0.521  -0.204   0.249
S3   90.663 101   0.760      0.889     0.914   -0.944  -0.724   0.432
S4  101.708 101   0.462      0.997     0.994    0.018  -0.014   0.276
S5  108.038 101   0.298      1.059     1.050    0.517   0.458   0.215
S6   96.788 101   0.600      0.949     0.916   -0.258  -0.538   0.347
S7  179.973 101   0.000      1.764     1.429    3.980   3.445  -0.041
S8  100.127 101   0.506      0.982     0.994   -0.032   0.030   0.106
S9   75.659 101   0.972      0.742     0.778   -1.618  -1.389   0.529
S10  82.940 101   0.905      0.813     0.800   -1.438  -1.809   0.530
S11  73.544 101   0.982      0.721     0.724   -2.576  -2.589   0.665
S12  88.530 101   0.808      0.868     0.877   -0.892  -0.955   0.422
\end{verbatim}

Analisis kesesuaian item awal menunjukkan bahwa 11 dari 12 item yang
dikalibrasi memiliki integritas struktural yang sangat baik, dengan
nilai yang berada dengan aman dalam batas kesesuaian praktis konservatif
antara 0,7 hingga 1,3. Hal ini mengindikasikan bahwa respons yang
diamati terhadap item-item tersebut sangat sesuai dengan ekspektasi
probabilistik model Rasch, yang mengonfirmasi konstruk pengukuran yang
bersih dan unidimensi.

Sebuah anomali psikometrik yang signifikan teridentifikasi pada Item S7.
Item ini menunjukkan profil ketidaksesuaian yang ekstrem (underfit),
dengan nilai Outfit MSQ sebesar 1,764, disertai indeks diskriminasi
negatif sebesar -0,041.

Item S7 menunjukkan ketidaksesuaian yang cukup kuat terhadap ekspektasi
model Rasch, sebagaimana tercermin dari nilai Outfit MSQ yang tinggi.
Nilai indeks diskriminasi negatif mengindikasikan bahwa pola respons
peserta tidak mengikuti arah yang diharapkan. Peserta dengan kemampuan
lebih rendah cenderung lebih sering memperoleh kategori jawaban yang
lebih tinggi atau benar, sedangkan peserta dengan kemampuan lebih tinggi
justru menunjukkan kecenderungan sebaliknya.

\paragraph{3.2 Separation Reliability}\label{separation-reliability}

Untuk menutupi keterbatasan prosedur CML yang sensitif terhadap sel
kosong pada uji invarian, presisi dan keandalan instrumen diverifikasi
menggunakan indeks \textbf{Reliability Separation}. Uji ini mengevaluasi
seberapa jauh instrumen memiliki daya pembeda yang andal untuk memetakan
gradasi kemampuan responden maupun hierarki kesulitan butir soal secara
konsisten menggunakan matriks data utuh (\(N = 102\)).

\begin{verbatim}
Rasch Reliability Metrics:
\end{verbatim}

\begin{verbatim}
Person Reliability (Separation):  0.62 
\end{verbatim}

\begin{verbatim}
Item Reliability (Separation)  :  0.905 
\end{verbatim}

\begin{itemize}
\tightlist
\item
  \textbf{Person Separation Reliability:} Mencatatkan koefisien sebesar
  \(0.62\). Nilai pada kategori sedang ini merefleksikan bahwa varians
  kemampuan teoretis di antara sampel acak berdistribusi cukup homogen.
\end{itemize}

\begin{itemize}
\tightlist
\item
  \textbf{Item Separation Reliability (Parameter Varians):} Mencatatkan
  koefisien yang sangat kuat, yakni \(0.905\). Angka ini membuktikan
  bahwa sebaran tingkat kesulitan ke-12 butir soal terkalibrasi dengan
  presisi sangat tinggi. Jarak antar-ambang kesulitan butir soal
  (\emph{logit scale}) sangat stabil dan andal, yang memastikan bahwa
  hierarki kesukaran soal akan mereplikasi pola identik apabila diujikan
  pada kelompok populasi lain yang setara.
\end{itemize}

\subsubsection{4. Tingkat Kesukaran Soal}\label{tingkat-kesukaran-soal}

\begin{verbatim}

Design Matrix Block 1:
    Location Threshold 1 Threshold 2 Threshold 3 Threshold 4
S1   1.13996     2.14635    -0.61988    -1.13939     4.17276
S2  -0.26814    -0.72817    -1.61408    -0.21779     1.48748
S3   1.70287     0.44718     2.95857          NA          NA
S4  -0.08668    -2.91431     0.30488     0.52865     1.73406
S5   0.64538     0.04046     0.18169     0.74833     1.61105
S6  -0.02914    -0.54046    -0.52178    -0.98036     1.92603
S7   0.29946     3.84541    -2.83638     1.30301    -1.11422
S8   0.21198     1.26632    -3.50799     2.05951     1.03006
S9   0.14619     1.08858    -3.16922     0.56570     2.09971
S10 -0.03455     0.05970    -0.10812    -0.23645     0.14665
S11  0.20022    -0.42339    -0.18102     0.49899     0.90629
S12  0.19040    -1.16204     0.26086    -1.05675     2.71951
\end{verbatim}

Matriks di atas menampilkan \textbf{Location} (tingkat kesulitan
rata-rata butir soal) serta \textbf{Thresholds} (titik potong transisi
antar jenjang kategori skor/jawaban) untuk ke-12 butir soal.

\paragraph{\texorpdfstring{4.1 Kolom \emph{Location} (Tingkat Kesulitan
Butir)}{4.1 Kolom Location (Tingkat Kesulitan Butir)}}\label{kolom-location-tingkat-kesulitan-butir}

\begin{itemize}
\item
  Menunjukkan posisi absolut item pada kontinum skala logit.
\item
  \textbf{Item Paling Sulit:} \textbf{S3} (1.70287) berada di posisi
  tertinggi, yang berarti butir ini membutuhkan kapasitas laten
  tertinggi.
\item
  \textbf{Item Paling Mudah:} \textbf{S2} (-0.26814) berada di posisi
  terendah.
\end{itemize}

\paragraph{\texorpdfstring{4.2 Kolom \emph{Thresholds 1--4} (Ambang
Batas Transisi
Kategori)}{4.2 Kolom Thresholds 1--4 (Ambang Batas Transisi Kategori)}}\label{kolom-thresholds-14-ambang-batas-transisi-kategori}

\begin{itemize}
\item
  Menunjukkan titik logit (\(\theta\)) di mana probabilitas responden
  untuk mencapai jenjang skor yang lebih tinggi mulai setara dengan
  peluang berada pada jenjang di bawahnya.
\item
  Semakin rendah/negatif nilai \emph{threshold}, semakin mudah suatu
  jenjang skor dicapai oleh responden.
\end{itemize}

\paragraph{4.3 Diagnostik Anomali: Kasus S3 dan
S7}\label{diagnostik-anomali-kasus-s3-dan-s7}

Dua butir soal ini menunjukkan karakteristik struktural yang menyimpang
dari pola umum:

\begin{itemize}
\item
  \textbf{Anomali S3 (\emph{Missing Thresholds}):} Terdapat nilai
  \texttt{NA} (Not Available) pada kolom \emph{Threshold 3} dan
  \emph{Threshold 4}. Hal ini mengindikasikan bahwa butir \textbf{S3}
  memiliki struktur penskoran yang lebih ringkas dibandingkan butir
  lainnya (misalnya hanya memiliki 3 jenjang kategori respons, bukan 4
  atau 5).
\item
  \textbf{Kekacauan S7 (\emph{Disordered Thresholds}):} Perhatikan
  urutan \emph{threshold} pada \textbf{S7} yang melompat tidak
  beraturan:

  \(3.84 \rightarrow -2.83 \rightarrow 1.30 \rightarrow -1.11\)

  Kondisi ini disebut \textbf{disordered thresholds} (ambang batas
  terbalik). Secara psikometri, ini adalah bukti matematis mutlak
  mengapa S7 mengalami \emph{misfit} parah (Outfit MSQ = 1.764) dan
  memiliki daya beda negatif. Jenjang jawaban pada butir ini tidak
  berfungsi secara hierarkis; responden dengan kemampuan lebih rendah
  secara anomali diuntungkan atau lebih mudah mencapai kategori skor
  tertentu dibandingkan responden dengan kemampuan lebih tinggi.
\end{itemize}

\subsubsection{5. Rangkuman Uji Tingkat Kesukaran, Uji Validitas, PTMEA
dan Daya
Beda}\label{rangkuman-uji-tingkat-kesukaran-uji-validitas-ptmea-dan-daya-beda}

\begin{verbatim}
   Item Calibration Infit_MSQ Outfit_MSQ Infit_ZSTD Outfit_ZSTD   PTMEA Discrim
1    S1      1.1400    0.7849     0.7489    -2.2819     -1.9091  0.6064  0.6064
2    S2     -0.2681    0.9616     0.9148    -0.2035     -0.5206  0.2493  0.2493
3    S3      1.7029    0.9138     0.8888    -0.7237     -0.9437  0.4325  0.4325
4    S4     -0.0867    0.9936     0.9971    -0.0143      0.0181  0.2758  0.2758
5    S5      0.6454    1.0500     1.0592     0.4577      0.5171  0.2155  0.2155
6    S6     -0.0291    0.9156     0.9489    -0.5383     -0.2576  0.3469  0.3469
7    S7      0.2995    1.4291     1.7644     3.4452      3.9805 -0.0410 -0.0410
8    S8      0.2120    0.9942     0.9816     0.0298     -0.0323  0.1056  0.1056
9    S9      0.1462    0.7775     0.7418    -1.3892     -1.6176  0.5295  0.5295
10  S10     -0.0346    0.8002     0.8131    -1.8085     -1.4383  0.5304  0.5304
11  S11      0.2002    0.7244     0.7210    -2.5890     -2.5761  0.6653  0.6653
12  S12      0.1904    0.8775     0.8679    -0.9553     -0.8921  0.4219  0.4219
\end{verbatim}

\subsubsection{6. Visualisasi Data}\label{visualisasi-data}

\paragraph{6.1 Wright Map}\label{wright-map}

\begin{center}
\pandocbounded{\includegraphics[keepaspectratio]{plot/Wright_Map.png}}
\end{center}

Sumbu horizontal menunjukkan skala logit (\(\theta\) / \(\beta\)), di
mana nilai yang semakin ke kanan mencerminkan kemampuan siswa yang lebih
tinggi atau soal yang lebih sukar.

\paragraph{Panel Atas: Distribusi Kemampuan Responden (Sisi
Siswa)}\label{panel-atas-distribusi-kemampuan-responden-sisi-siswa}

Bagian atas menampilkan kurva kepadatan (\emph{density curve}) yang
merepresentasikan sebaran tingkat kemampuan laten (\(\theta\)) dari
\(102\) siswa sampel, yang dikelompokkan berdasarkan profil kemampuan:

\begin{itemize}
\item
  \textbf{Kelompok Kemampuan Rendah (Area Biru):} Berpusat di sekitar
  wilayah logit \(-0.4\), dengan puncak minor di \(-0.9\). Ini
  menunjukkan konsentrasi mayoritas siswa yang memiliki kapasitas
  kognitif lebih rendah dibandingkan tingkat kesulitan rata-rata soal.
\item
  \textbf{Kelompok Kemampuan Sedang (Area Kuning):} Tersebar di rentang
  logit \(-0.2\) hingga \(0.5\), menjembatani antara siswa berkemampuan
  rendah dan tinggi.
\item
  \textbf{Kelompok Kemampuan Tinggi (Area Merah):} Membentuk kurva
  landai di rentang logit \(0.5\) hingga \(1.5\). Ini menunjukkan
  sebagian kecil sampel yang memiliki kapasitas penalaran tingkat
  tinggi.
\end{itemize}

\paragraph{Panel Bawah: Kalibrasi Kesulitan Butir Soal (Sisi
Soal)}\label{panel-bawah-kalibrasi-kesulitan-butir-soal-sisi-soal}

Bagian bawah memetakan posisi ke-12 butir soal (S1 hingga S12)
berdasarkan tingkat kesulitannya (\(\beta\)). Butir-butir soal ini
diurutkan secara vertikal dan diletakkan pada titik logit yang sesuai:

\begin{itemize}
\item
  \textbf{Butir Soal Paling Sulit:} \textbf{S3} (berada paling kanan di
  sekitar logit \(+1.7\)). Soal ini menuntut kapasitas kemampuan siswa
  yang sangat tinggi untuk dapat diselesaikan dengan probabilitas
  keberhasilan \(0.5\).
\item
  \textbf{Butir Soal Menengah-Atas:} \textbf{S1} (logit \(+1.1\)) dan
  \textbf{S5} (logit \(+0.5\)).
\item
  \textbf{Butir Soal Menengah:} Kelompok soal \textbf{S6, S7, S8, S9,
  S11, S12, dan S4} mengelompok di sekitar wilayah logit \(0.0\) hingga
  \(0.3\). Klaster ini menunjukkan bahwa mayoritas soal rancangan pakar
  memiliki tingkat kesukaran yang setara dengan kemampuan siswa
  menengah.
\item
  \textbf{Butir Soal Paling Mudah:} \textbf{S2} (logit \(-0.3\)) dan
  \textbf{S10} (logit \(-0.1\)). Butir ini dapat dijangkau oleh sebagian
  besar siswa, termasuk mereka yang berada di kelompok kemampuan bawah.
\end{itemize}

\paragraph{\texorpdfstring{Evaluasi Penargetan Instrumen
(\emph{Targeting}) dan Garis
Referensi}{Evaluasi Penargetan Instrumen (Targeting) dan Garis Referensi}}\label{evaluasi-penargetan-instrumen-targeting-dan-garis-referensi}

Dua garis putus-putus vertikal dan metrik di bagian judul memberikan
konfirmasi statistik mengenai kualitas perancangan instrumen:

\begin{itemize}
\item
  \textbf{Garis Putus-Putus Biru (Mean Person Ability} \(\approx +0.22\)
  logits): Titik tengah dari keseluruhan kemampuan siswa.
\item
  \textbf{Garis Putus-Putus Merah (Mean Item Difficulty}
  \(\approx +0.25\) logits): Titik tengah dari tingkat kesulitan butir
  soal.
\item
  \textbf{Mean Gap (}\(0.15\) logits): Jarak antara rata-rata kemampuan
  siswa dan rata-rata kesulitan soal sangat tipis. Ini berarti instrumen
  secara teoretis diletakkan pada wilayah yang tepat untuk mengukur
  populasi sampel ini.
\item
  \textbf{IPOI =} \(0.942\) (Excellent Targeting): Nilai
  \emph{Item-Person Separation Reliability} atau indeks penargetan yang
  mendekati \(1.0\) menunjukkan bahwa instrumen ini memiliki rentang
  kesukaran yang sangat ideal untuk memetakan disparitas kemampuan siswa
  yang diuji. Tidak ada jurang pemisah yang ekstrem di mana soal terlalu
  mustahil dikerjakan atau terlalu sepele untuk seluruh populasi.
\end{itemize}

Wright Map ini menceritakan dinamika antara pembuat soal dan realitas
siswa di lapangan:

\begin{enumerate}
\def\labelenumi{\arabic{enumi}.}
\item
  \textbf{Kesenjangan Puncak Kemampuan dan Kesukaran:} Perhatikan bahwa
  puncak kurva biru (mayoritas siswa) berada di sebelah kiri (sekitar
  \(-0.4\)), sedangkan sebaran soal (titik merah) lebih banyak menumpuk
  di sebelah kanan (sekitar \(0.0\) hingga \(+0.3\)). Hal ini
  mengindikasikan bahwa meskipun instrumen tertarget dengan baik secara
  keseluruhan, instrumen ini \textbf{relatif sedikit lebih
  menantang/sukar} bagi konsentrasi terbesar siswa (kelompok kemampuan
  rendah).
\item
  \textbf{Daya Diskriminasi Soal S3:} Butir \textbf{S3} yang terisolasi
  di sisi paling kanan (\(+1.7\)) berfungsi secara spesifik sebagai
  pembeda mutlak (\emph{discriminator}) untuk menyaring siswa-siswa
  dengan tingkat penguasaan kognitif tertinggi di kelas populasi.
\item
  \textbf{Penyebaran Soal:} Sebaran titik merah pada sumbu logit
  menunjukkan bahwa instrumen memiliki gradasi atau tangga kesulitan
  yang bervariasi, tidak menumpuk di satu titik saja, sehingga instrumen
  ini informatif dalam memetakan gradasi kapasitas kognitif responden.
\end{enumerate}

\paragraph{6.2 Item Fit}\label{item-fit-1}

\begin{center}
\pandocbounded{\includegraphics[keepaspectratio]{plot/item_fit.png}}
\end{center}

Plot \textbf{Item Difficulty and Outfit Statistics} adalah representasi
visual untuk mengaudit validitas konstruk dari 12 butir soal instrumen.
Plot ini memetakan posisi setiap soal berdasarkan tingkat kesulitannya
sekaligus mendeteksi apakah soal tersebut berfungsi sebagaimana mestinya
secara statistik.

\paragraph{Membaca Sumbu dan Garis
Referensi}\label{membaca-sumbu-dan-garis-referensi}

\begin{itemize}
\item
  \textbf{Sumbu X (\emph{Item Difficulty} / Skala Logit):} Menunjukkan
  tingkat kesulitan butir soal.

  \begin{itemize}
  \item
    Semakin ke kanan posisi titik, semakin sukar soal tersebut.
  \item
    Contoh: \textbf{S3} adalah soal paling sukar (berada di dekat logit
    +1.7), sedangkan \textbf{S2} adalah soal yang paling mudah (berada
    di dekat logit -0.3).
  \end{itemize}
\item
  \textbf{Sumbu Y (\emph{Outfit Mean Square} / MSQ):} Menunjukkan
  tingkat penyimpangan atau \emph{noise} (keacakan) respons siswa.

  \begin{itemize}
  \item
    \textbf{Garis putus-putus horizontal di 1.0:} Nilai ideal teoretis.
    Semakin dekat titik ke garis ini, semakin sempurna pola respons
    siswa terhadap model Rasch.
  \item
    Secara psikometri, batas toleransi penerimaan (\emph{practical fit})
    untuk instrumen non-standar adalah \textbf{0.7 hingga 1.3}. Titik
    yang berada di atas 1.3 mengindikasikan adanya keacakan, sementara
    di bawah 0.7 mengindikasikan kebergantungan berlebih
    (\emph{redundancy}).
  \end{itemize}
\end{itemize}

\paragraph{\texorpdfstring{Klaster Validitas: Kelompok \emph{Fit}
(Lingkaran
Abu-abu)}{Klaster Validitas: Kelompok Fit (Lingkaran Abu-abu)}}\label{klaster-validitas-kelompok-fit-lingkaran-abu-abu}

Sebelas dari dua belas butir soal---meliputi \textbf{S1, S2, S3, S4, S5,
S6, S8, S9, S10, S11, dan S12}---digambarkan sebagai lingkaran abu-abu
dan diklasifikasikan sebagai \textbf{Fit}.

\begin{itemize}
\item
  Kesemua titik ini berada di dalam jalur aman, yakni di bawah garis
  batas \emph{Outfit MSQ} 1.0 (bahkan semuanya berada di rentang 0.7
  hingga 1.1).
\item
  Hal ini membuktikan bahwa varians respons siswa pada 11 soal tersebut
  dapat diprediksi dengan sangat baik oleh probabilitas model Rasch.
  Data empiris ini mengonfirmasi \textbf{validitas konstruk} instrumen
  Anda secara sangat kuat.
\end{itemize}

\paragraph{\texorpdfstring{Anomali Diagnostik: Butir \emph{Misfit}
(Segitiga Hitam
S7)}{Anomali Diagnostik: Butir Misfit (Segitiga Hitam S7)}}\label{anomali-diagnostik-butir-misfit-segitiga-hitam-s7}

Secara visual maupun numerik, \textbf{Butir Soal S7} (segitiga hitam)
memisahkan diri secara ekstrem dari klaster lainnya dan melompat jauh ke
atas, menyentuh nilai \emph{Outfit MSQ} hampir \textbf{1.8}.

\begin{itemize}
\item
  \textbf{Deviasi Respons:} Nilai \emph{Outfit} sebesar 1.764
  mengindikasikan bahwa terdapat 76.4\% keacakan atau \emph{noise} dalam
  data respons siswa yang tidak dapat dijelaskan oleh model.
\item
  Seperti yang telah dikalkulasikan sebelumnya, anomali ini sejalan
  dengan nilai daya beda empiris (korelasi poin-biserial laten) yang
  negatif (\textbf{-0.041}). Secara substansial, ini adalah bukti visual
  bahwa S7 gagal membedakan kemampuan siswa. Terdapat indikasi kuat
  bahwa pengecoh (\emph{distractor}), kompleksitas kalimat, atau muatan
  kontekstual pada soal S7 menyesatkan siswa yang berkemampuan tinggi
  atau sebaliknya, memberikan keuntungan tebakan bagi siswa berkemampuan
  rendah.
\end{itemize}




\end{document}
